\documentclass[10pt,a4paper]{article}
\usepackage[margin=1.7cm]{geometry}
\usepackage{enumitem}\usepackage{titlesec}\usepackage[hidelinks]{hyperref}\usepackage{parskip}
\titleformat{\section}{\normalfont\large\bfseries}{\thesection}{0.6em}{}
\titlespacing*{\section}{0pt}{6pt}{2pt}
\setlist[itemize]{topsep=2pt,itemsep=1pt,parsep=0pt,leftmargin=1.3em}
\pagestyle{empty}
\begin{document}
\begin{center}
{\Large\bfseries FPGA Image Processing: VGA Display and a 3x3 Median Filter}\\[2pt]
{\large Verilog RTL designs on a Xilinx FPGA}\\[4pt]
{\small FPGA laboratory course \textendash{} B.Sc. Electrical Engineering, Isfahan University of Technology}\\
{\small Author: Kiavash Rabbani Soudabi (student ID 9821663) \textendash{} term 4012 (Fall 2022) \textendash{} individual work}
\end{center}

\section*{Purpose}
Build two hardware image-processing designs in Verilog and verify them through simulation and on the lab board:
(1) display a 224x224 image stored in FPGA block memory on a VGA monitor, and (2) filter a stream of image pixels in
hardware with a 3x3 median filter, which is the standard pre-processing step before edge detection such as Sobel.

\section*{Design 1 \textendash{} VGA image display}
The top level instantiates a clock-wizard IP core that derives the 25 MHz and 40 MHz pixel clocks from the 24 MHz
board oscillator, an image ROM initialised from a Xilinx \texttt{.coe} file, and a pixel generator that walks the
stored image and drives the RGB outputs. Dedicated synchronisation modules generate the horizontal and vertical
timing, and a UCF file constrains the pins and the board clock. Test patterns of 224x224 pixels (3-bit and 8-bit)
are kept in the repository as ROM initialisation data.

\section*{Design 2 \textendash{} 3x3 median filter}
The second design streams pixels from an 8-bit block-memory ROM (depth 50176 = 224x224) and computes the median of
every 3x3 neighbourhood on the fly:
\begin{itemize}
  \item nine registers hold the current window; the column registers are shifted one position per clock after every
        third clock, so a complete window is always available without a frame buffer;
  \item the median is extracted with a compare-and-swap network over those registers, and the filtered pixel is
        written to the output stream together with an \texttt{op\_finished} flag;
  \item a 16-bit address counter with a comparison constant marks the end of each three-line block, which keeps the
        window aligned with the image rows.
\end{itemize}

\section*{Verification}
The filter was verified with Verilog testbenches before programming the device; the captured simulation waveform is
included in the repository. The ROM images (a noisy photograph used as the median-filter input and a grayscale scene
used for edge detection) were decoded from the committed \texttt{.coe} files and are shown in the README, so the
processed content can be inspected without hardware.

\section*{Tools}
Verilog-2001, Xilinx ISE (synthesis, UCF constraints, implementation) with ISim/ModelSim for simulation, Xilinx IP
cores (clock wizard, block memory generator), VGA display.

\section*{Repository contents}
\texttt{exp6-vga/} (VGA design, pre-lab code, UCF and IP sources) $\cdot$ \texttt{exp7-median-filter/} (median filter,
testbenches, ROM initialisation) $\cdot$ \texttt{data/} (the 224x224 test images in \texttt{.coe} format) $\cdot$
\texttt{reports/} (submitted pre-report) $\cdot$ \texttt{docs/} (this overview, simulation waveform and decoded test
images).

\section*{Skills demonstrated}
Verilog RTL design; FPGA design flow in Xilinx ISE; VGA timing and synchronisation; clock generation with IP cores;
block-memory based image storage; streaming image processing in hardware; 3x3 windowing by column shifting; median
filter implementation as a compare-and-swap network; testbench-based verification of RTL.
\end{document}
